% Bewertungspaket zum Gesamttest «Trigonometrische Berechnungen» — Quelle des PDFs
% (python3 scripts/build-lp-pdf.py trigonometrische-berechnungen). Ganze Punkte je Teilschritt; Werte mit python3
% nachgerechnet (07.10.2026, Folgefehler-Fälle eingeschlossen). Aufbau wie planimetrie/bewertungspaket.tex.
\documentclass[11pt]{article}
\usepackage{../lp-druck}
\renewcommand{\lpname}{Trigonometrische Berechnungen}
\renewcommand{\lpbereich}{Grundlagenfach}
\renewcommand{\lpteilgebiet}{5.3}
\renewcommand{\lpfuss}{Bewertungspaket}
\newcolumntype{P}{>{\centering\arraybackslash}p{10mm}}
\newenvironment{raster}{\par\smallskip\noindent\tabularx{\linewidth}{@{}>{\raggedright\arraybackslash}p{38mm}P X@{}}\toprule
  \small\textsc{Teilschritt} & \small\textsc{P} & \small\textsc{Musterlösung}\\\midrule}{\bottomrule\endtabularx\par}
\newcommand{\fehler}[1]{\par\smallskip{\small\textcolor{lprot}{\bfseries Typische Fehler:} #1}}
\newcommand{\E}{\,{\footnotesize\bfseries\color{lpgruen}(E)}}
\newcommand{\A}{\,{\footnotesize\bfseries\color{lpblau}(A)}}
\begin{document}
\lpkopf{Bewertungspaket zum Gesamttest}{}

\begin{lpachtung}{Erst lösen, dann öffnen}
Dieses Paket enthält die Musterlösung. Wer es vor dem Lösen liest, misst am Ende nur, wie gut das Abschreiben gelingt.
\end{lpachtung}

\section*{1 · So gehst du vor}
\begin{enumerate}[leftmargin=*,itemsep=2pt]
\item \textbf{Gesamttest lösen} — vollständig, mit Skizze und Rechenweg, ohne dieses Paket; der Taschenrechner ist erlaubt (Gradmodus).
\item \textbf{Fotografieren oder scannen}, gut lesbar, eine Seite pro Bild. \textbf{Kein Name} auf den Bildern.
  Handyfotos tragen oft unsichtbar Standort, Zeit und Gerät mit (Metadaten). Lade darum lieber einen Scan oder ein
  Bildschirmfoto des Fotos hoch, oder entferne vorher die Standortangaben.
\item \textbf{Bewerten lassen.} Ob du dafür eine KI nutzt und welche, entscheidest du selbst und in eigener Verantwortung —
  je nachdem, was dir aufgrund deines Alters und deiner persönlichen Situation erlaubt ist (Altersgrenzen und
  Nutzungsbedingungen des Anbieters, allenfalls Einverständnis deiner Eltern). Lade nur deine Lösung und dieses PDF hoch,
  keine persönlichen Daten, und füge den Auftrag aus Abschnitt~2 ein. Ohne KI kannst du dich mit Abschnitt~3 selbst bewerten.
\item \textbf{Rückmeldung prüfen.} Stimmt, was die KI aus deiner Handschrift gelesen hat? Eine KI kann sich irren —
  im Zweifel gilt das Raster in Abschnitt~3.
\item \textbf{Gezielt wiederholen} nach Abschnitt~4.
\end{enumerate}

\needspace{8\baselineskip}
\section*{2 · Auftrag an die KI}
\begin{tcolorbox}[breakable,colback=lplinie!25,colframe=lplinie,boxrule=0.5pt,arc=1mm,fontupper=\small]
Du bist Korrektorin oder Korrektor für Mathematik an einer Schweizer Berufsmaturitätsschule. Im Anhang findest du (1)~das Bewertungspaket zum Gesamttest «Trigonometrische Berechnungen» mit Musterlösung und Punkteraster und (2)~Fotos meiner handschriftlichen Lösung.\par\medskip
Geh so vor:\par
1. Schreib zu jeder Aufgabe G1 bis G7 zuerst ab, was du in meiner Lösung liest — Rechenweg und Ergebnis. Was du nicht sicher lesen kannst, markierst du mit [unsicher]. Nicht raten.\par
2. Vergib die Punkte genau nach dem Raster im Bewertungspaket (Abschnitt 3), ganze Punkte je Zeile. Ziehe einen Fehler nur einmal ab: Rechne ich mit einem falschen Zwischenergebnis richtig weiter, gibt es die Folgepunkte. Ausnahme: Zeilen mit \textbf{(E)} gibt es nur für das richtige Ergebnis, nie als Folgepunkt. Die Zeilen «Typische Fehler» sagen, welche Rasterzeile bei diesem Fehler 0 Punkte gibt — sie sind kein zusätzlicher Abzug.\par
3. Andere richtige Lösungswege zählen voll, ebenso gleichwertige Schreibweisen: Dezimalpunkt oder Dezimalkomma, exakte Werte wie $\sqrt{35.24}$ statt Dezimalzahlen, anders, aber vernünftig gerundete Ergebnisse (Abweichung höchstens $\pm 0.01$ in der letzten Stelle oder durch Runden von Zwischenergebnissen erklärbar). \textbf{Einheiten:} Fehlt bei einem Endergebnis die Einheit oder ist sie falsch ($\text{cm}$ statt $\text{cm}^2$), gibt die (E)-Zeile dieses Ergebnisses 0 Punkte — aber nur beim ersten Mal im ganzen Gesamttest; weitere Einheitenfehler kosten nichts mehr. Zeilen mit \textbf{(A)} verlangen nur Erkennen oder Ablesen, diese Punkte gibt es auch ohne Rechenweg. Zeilen mit \textbf{(E)} sind Ergebnispunkte: Das Ergebnis muss stimmen \emph{und} aus einem erkennbaren Weg hervorgehen. Alle übrigen Zeilen bewerten den Weg selbst.\par
4. Begründe jeden Abzug in einem Satz und nenne die Fehlerart (zum Beispiel «Gegen- und Ankathete vertauscht» oder «falsches Paar im Sinussatz»). Schreib die Musterlösung nicht ab — zeig mir, wo mein Fehler liegt.\par
5. Gib zum Schluss eine Tabelle «Aufgabe | Punkte | Begründung», die Gesamtpunktzahl (von 25) und — nach der Selbsteinschätzung in Abschnitt 4 — welches Kapitel des Leitprogramms ich wiederholen soll. Keine Note.\par
6. Fehlt ein Foto oder ist eine Aufgabe unlesbar, sag es, statt sie zu bewerten.
\end{tcolorbox}

\needspace{8\baselineskip}
\section*{3 · Musterlösung und Punkteraster}
\textbf{Allgemein:} Ganze Punkte je Zeile. Ein Fehler wird nur einmal abgezogen (Folgefehler: wer mit einem falschen
Zwischenergebnis richtig weiterrechnet, bekommt die Wegpunkte der folgenden Zeilen, aber keine (E)-Punkte). In der Spalte P:
\textbf{(A)}~= erkennen oder ablesen, zählt auch ohne Rechenweg · \textbf{(E)}~= Ergebnispunkt, Ergebnis richtig
\emph{und} Weg erkennbar · ohne Zeichen = die Zeile bewertet den Weg. Taschenrechner erlaubt; Längen und Winkel auf zwei
Dezimalen, Toleranz $\pm 0.01$ (oder durch gerundete Zwischenwerte erklärbar). \textbf{Einheiten:} Der erste Einheitenfehler im
ganzen Test kostet die (E)-Zeile dieses Ergebnisses, weitere nichts. \textbf{Gradmodus:} Ein Ergebnis aus dem Bogenmass (RAD)
ist falsch; die Wegzeilen zählen, die (E)-Zeilen nicht.
«Typische Fehler» nennt, welche Zeile 0 Punkte gibt, und die Punkte, die bleiben.
\textbf{Teilrichtige Zeilen:} Verlangt eine Zeile mehrere Angaben und stimmt nur ein Teil, gibt die Zeile 0 Punkte.

\begin{lpaufgabe}{G1}{Seiten und Winkel}{4 P}
\begin{raster}
(a) Seiten benennen & 1\A & Von $P$ aus: Gegenkathete $\overline{QR}$, Ankathete $\overline{PQ}$ ($\overline{PR}$ ist die Hypotenuse). Beide nötig.\\
(b) $\overline{PQ}$ & 1\E & $\overline{PQ} = 11 \cdot \cos 36^\circ \approx 8.90\,\text{cm}$\\
(b) $\overline{QR}$ & 1\E & $\overline{QR} = 11 \cdot \sin 36^\circ \approx 6.47\,\text{cm}$\\
(c) Winkel bei $R$ & 1\E & $90^\circ - 36^\circ = 54^\circ$\\
\end{raster}
\fehler{Gegen- und Ankathete vertauscht ($\overline{PQ} \approx 6.47$, $\overline{QR} \approx 8.90$): (a) und beide (b)-Zeilen 0 $\to$ 1 von 4\,P.
$11 : \cos 36^\circ \approx 13.60$ (geteilt statt multipliziert, länger als die Hypotenuse): diese (b)-Zeile 0.
(c) $180^\circ - 36^\circ = 144^\circ$ (rechten Winkel vergessen): (c) 0.}
\end{lpaufgabe}

\begin{lpaufgabe}{G2}{Satteldach}{3 P}
\begin{raster}
(a) Ansatz & 1 & Halbes Dach: rechtwinkliges Dreieck mit Ankathete $5\,\text{m}$ und Gegenkathete $3.2\,\text{m}$; $\tan x = \tfrac{3.2}{5}$\\
(a) Winkel & 1\E & $x = \arctan 0.64 \approx 32.62^\circ$\\
(b) Sparren & 1\E & $\sqrt{5^2 + 3.2^2} = \sqrt{35.24} \approx 5.94\,\text{m}$ (gleichwertig: $\tfrac{5}{\cos x}$ oder $\tfrac{3.2}{\sin x}$)\\
\end{raster}
\fehler{Die ganze Breite $10\,\text{m}$ als Ankathete ($x \approx 17.74^\circ$, Sparren $\approx 10.50\,\text{m}$): Ansatz und beide (E)-Zeilen 0 $\to$ 0 von 3\,P.
$x = 0.64$ oder $x = 64^\circ$ (Verhältnis statt Winkel): Winkelzeile 0, der Rest zählt.}
\end{lpaufgabe}

\begin{lpaufgabe}{G3}{Zwei Häuser vom Turm aus}{4 P}
\begin{raster}
Skizze & 1 & Turm $32\,\text{m}$ senkrecht, Waagrechte durch die Plattform, beide Sehstrahlen; die Tiefenwinkel gegen die Waagrechte und als gleich grosse Höhenwinkel an den Häusern (Wechselwinkel).\\
näheres Haus & 1\E & $\tan 24^\circ = \tfrac{32}{d_1}$, $d_1 = \tfrac{32}{\tan 24^\circ} \approx 71.87\,\text{m}$\\
ferneres Haus & 1\E & $d_2 = \tfrac{32}{\tan 15^\circ} \approx 119.43\,\text{m}$\\
Abstand & 1\E & $d_2 - d_1 \approx 47.55\,\text{m}$ (mit den gerundeten Werten $47.56\,\text{m}$: zählt)\\
\end{raster}
\fehler{$32 \cdot \tan 24^\circ \approx 14.25$ und $32 \cdot \tan 15^\circ \approx 8.57$ (Gegen- und Ankathete vertauscht; der nähere Abstand wäre der grössere): beide Abstandszeilen und die letzte Zeile 0 $\to$ höchstens 1 von 4\,P.
Tiefenwinkel gegen die Senkrechte gemessen ($66^\circ$ und $75^\circ$ als Höhenwinkel an den Häusern): Skizzenzeile 0, die drei (E)-Zeilen 0 $\to$ 0 von 4\,P.}
\end{lpaufgabe}

\begin{lpaufgabe}{G4}{Waldbrand}{4 P}
\begin{raster}
Dritter Winkel & 1 & $\angle F = 180^\circ - 48^\circ - 64^\circ = 68^\circ$\\
(a) Ansatz & 1 & Sinussatz mit dem Paar $\overline{AB}$, $\angle F$: $\tfrac{\overline{AF}}{\sin 64^\circ} = \tfrac{6.5}{\sin 68^\circ}$ ($\overline{AF}$ liegt dem Winkel bei $B$ gegenüber)\\
(a) $\overline{AF}$ & 1\E & $\overline{AF} \approx 6.30\,\text{km}$\\
(b) Abstand & 1\E & Rechtwinkliges Dreieck mit Hypotenuse $\overline{AF}$: $h = \overline{AF} \cdot \sin 48^\circ \approx 4.68\,\text{km}$ (gleichwertig: $\overline{BF} \cdot \sin 64^\circ$ mit $\overline{BF} \approx 5.21\,\text{km}$)\\
\end{raster}
\fehler{Falsches Paar, $\overline{AF} = \tfrac{6.5 \cdot \sin 48^\circ}{\sin 68^\circ} \approx 5.21$ (das ist $\overline{BF}$): Ansatz und (a)-Ergebnis 0; (b) mit diesem Wert ($\approx 3.87$) 0 $\to$ 1 von 4\,P.
Dritten Winkel vergessen und mit $\sin 90^\circ$ gerechnet: Winkel-, Ansatz- und Ergebniszeile 0.
(b) $\overline{AF} \cdot \cos 48^\circ$ (Abstand am Boden statt Höhe): (b) 0.}
\end{lpaufgabe}

\begin{lpaufgabe}{G5}{Zwei mögliche Dreiecke}{3 P}
\begin{raster}
(a) Anzahl & 1 & $h = 10 \cdot \sin 42^\circ \approx 6.69\,\text{cm}$; wegen $h < a < c$ ($6.69 < 8 < 10$) schneidet der Kreis um $B$ den Schenkel zweimal: zwei Dreiecke. Ohne Vergleich mit $h$: 0.\\
(b) $\gamma_2$ & 1\E & $\sin\gamma = \tfrac{10 \cdot \sin 42^\circ}{8} \approx 0.8364$, $\gamma_1 \approx 56.76^\circ$, $\gamma_2 = 180^\circ - \gamma_1 \approx 123.24^\circ$\\
(b) $\beta$ & 1\E & $\beta = 180^\circ - 42^\circ - 123.24^\circ \approx 14.76^\circ$\\
\end{raster}
\fehler{Nur den Rechnerwert $\gamma \approx 56.76^\circ$ und $\beta \approx 81.24^\circ$ (das spitze Dreieck): beide (b)-Zeilen 0 $\to$ höchstens 1 von 3\,P.
$\sin\gamma = \tfrac{8 \cdot \sin 42^\circ}{10}$ (Paar vertauscht): (b) 0.}
\end{lpaufgabe}

\begin{lpaufgabe}{G6}{Grundstück}{4 P}
\begin{raster}
(a) Ansatz & 1 & Cosinussatz (zwei Seiten und der eingeschlossene Winkel): $s^2 = 60^2 + 85^2 - 2 \cdot 60 \cdot 85 \cdot \cos 125^\circ$\\
(a) Seite & 1\E & $s^2 \approx 16\,675.48$, $s \approx 129.13\,\text{m}$\\
(b) Ansatz & 1 & $A = \tfrac{60 \cdot 85}{2} \cdot \sin 125^\circ$\\
(b) Fläche & 1\E & $A \approx 2088.84\,\text{m}^2$\\
\end{raster}
\fehler{Korrekturglied vergessen ($\sqrt{60^2 + 85^2} \approx 104.04$): (a) Ansatz und Ergebnis 0.
Vorzeichen falsch ($+2 \cdot 60 \cdot 85 \cdot \cos 125^\circ$, ergibt $\approx 70.53$ — kürzer als eine gegebene Seite trotz stumpfem Winkel): (a)-Ergebnis 0.
Fläche mit $\cos$ statt $\sin$ (negativ): (b) Ansatz und Ergebnis 0. $\tfrac12$ vergessen ($\approx 4177.68$): (b)-Ergebnis 0.}
\end{lpaufgabe}

\begin{lpaufgabe}{G7}{Eine fremde Lösung}{3 P}
\begin{raster}
Urteil & 1 & Nein: Ein negativer Cosinus gehört zu einem stumpfen Winkel; Lias Rechnung stimmt.\\
Winkel & 1\E & $\gamma = \arccos(-0.65) \approx 130.54^\circ$\\
Begründung & 1 & $c^2 = 100 > a^2 + b^2 = 61$: Die Seite gegenüber $\gamma$ ist länger als im rechtwinkligen Fall, also ist $\gamma$ stumpf (gleichwertig: das Korrekturglied $-2ab\cos\gamma$ ist positiv).\\
\end{raster}
\fehler{$\gamma = \arccos(0.65) \approx 49.46^\circ$ (Vorzeichen weggelassen): Winkelzeile 0, Urteil 0.
«Ja, verrechnet»: Urteil 0; die übrigen Zeilen zählen, wenn sie stimmen.}
\end{lpaufgabe}

\needspace{14\baselineskip}
\section*{4 · Selbsteinschätzung}
\begin{tabularx}{\linewidth}{@{}l X@{}}\toprule
22 – 25 P & Die geprüften Teile sitzen. Wo du Punkte verloren hast: das Kapitel dieser Aufgabe nochmals (Zuordnung unten).\\
17 – 21 P & Den schwächsten Teil nochmals: Tüfteln und Übungen des Kapitels, in dem du die meisten Punkte verloren hast.\\
11 – 16 P & Zurück zu den Kapiteln aller Aufgaben, in denen du Punkte verloren hast.\\
0 – 10 P & Zurück zu Kapitel 1 und von dort der Reihe nach weiter.\\\midrule
\multicolumn{2}{@{}p{\linewidth}@{}}{\small\textbf{Aufgabe $\to$ Kapitel:} G1 $\to$ Kapitel 1 (Sinus, Cosinus und Tangens); G2 $\to$ Kapitel 2 (Winkel berechnen); G3 $\to$ Kapitel 3 (Höhen und Distanzen); G4, G5 $\to$ Kapitel 4 (Sinussatz); G6, G7 $\to$ Kapitel 5 (Cosinussatz und Fläche)}\\\bottomrule
\end{tabularx}
\end{document}
